\textbf{Phần 1. Trắc nghiệm lựa chọn}
\setcounterpageref{ex}{0}

\begin{ex} %[2D8N40]
Kết quả của phép tính $5 \times 3$ là
\choice
{\True $15$}
{$10$}
{$20$}
{$25$}
\loigiai{
$5 \times 3 = 15$.\\
\textbf{Đáp án của câu hỏi là A}
}
\end{ex}

\begin{ex} %[2D8N40]
Kết quả của phép tính $5 \times 6$ là
\choice
{$25$}
{$35$}
{\True $30$}
{$40$}
\loigiai{
$5 \times 6 = 30$.\\
\textbf{Đáp án của câu hỏi là C}
}
\end{ex}

\begin{ex} %[2D8H40]
Phép tính nào sau đây có kết quả bằng 40?
\choice
{$5 \times 7$}
{\True $5 \times 8$}
{$5 \times 9$}
{$5 \times 6$}
\loigiai{
$5 \times 8 = 40$.\\
\textbf{Đáp án của câu hỏi là B}
}
\end{ex}

\begin{ex} %[2D8H40]
Tất cả các tích trong bảng nhân 5 đều có chữ số tận cùng là
\choice
{$0$ hoặc $2$}
{$5$ hoặc $9$}
{Chỉ là chữ số $5$}
{\True $0$ hoặc $5$}
\loigiai{
Các tích bảng nhân 5 luôn tận cùng là 0 hoặc 5.\\
\textbf{Đáp án của câu hỏi là D}
}
\end{ex}

\begin{ex} %[2D8H40]
Trong bảng nhân 5, hai tích liên tiếp hơn kém nhau bao nhiêu đơn vị?
\choice
{$2$ đơn vị}
{$10$ đơn vị}
{\True $5$ đơn vị}
{$1$ đơn vị}
\loigiai{
Hơn kém nhau 5 đơn vị.\\
\textbf{Đáp án của câu hỏi là C}
}
\end{ex}

\begin{ex} %[2D8H40]
Điền số thích hợp: $5 \times \dots = 45$.
\choice
{$8$}
{\True $9$}
{$7$}
{$10$}
\loigiai{
$5 \times 9 = 45$.\\
\textbf{Đáp án của câu hỏi là B}
}
\end{ex}

\begin{ex} %[2D8V40]
Mỗi bàn tay có 5 ngón tay. Hỏi 6 bàn tay có tất cả bao nhiêu ngón tay?
\choice
{\True $30$ ngón}
{$25$ ngón}
{$35$ ngón}
{$40$ ngón}
\loigiai{
$5 \times 6 = 30$ ngón tay.\\
\textbf{Đáp án của câu hỏi là A}
}
\end{ex}

\begin{ex} %[2D8V40]
Mỗi túi có 5 quả cam. Hỏi 8 túi như thế có tất cả bao nhiêu quả cam?
\choice
{$35$ quả}
{\True $40$ quả}
{$45$ quả}
{$50$ quả}
\loigiai{
$5 \times 8 = 40$ quả cam.\\
\textbf{Đáp án của câu hỏi là B}
}
\end{ex}

\begin{ex} %[2D8V40]
Tính giá trị biểu thức: $5 \times 7 + 15$.
\choice
{$45$}
{$55$}
{$60$}
{\True $50$}
\loigiai{
$35 + 15 = 50$.\\
\textbf{Đáp án của câu hỏi là D}
}
\end{ex}

\begin{ex} %[2D8V40]
Điền dấu thích hợp: $5 \times 4 \dots 2 \times 10$.
\choice
{$>$}
{$<$}
{\True $=$}
{Không so sánh được}
\loigiai{
$5 \times 4 = 20$ và $2 \times 10 = 20$. Vậy bằng nhau.\\
\textbf{Đáp án của câu hỏi là C}
}
\end{ex}

\begin{ex} %[2D8V40]
Một ngôi sao 5 cánh. Hỏi 9 ngôi sao như thế có tất cả bao nhiêu cánh?
\choice
{\True $45$ cánh}
{$40$ cánh}
{$50$ cánh}
{$35$ cánh}
\loigiai{
$5 \times 9 = 45$ cánh.\\
\textbf{Đáp án của câu hỏi là A}
}
\end{ex}

\begin{ex} %[2D8C40]
Có 7 can dầu loại $5\text{ l}$. Người ta đã dùng hết $18\text{ l}$ dầu. Hỏi còn lại bao nhiêu lít dầu?
\choice
{$16\text{ l}$}
{$18\text{ l}$}
{$19\text{ l}$}
{\True $17\text{ l}$}
\loigiai{
Tổng dầu: $5 \times 7 = 35\text{ l}$.\\
Còn lại: $35 - 18 = 17\text{ l}$.\\
\textbf{Đáp án của câu hỏi là D}
}
\end{ex}

\textbf{Phần 2. Trắc nghiệm đúng sai}
\setcounterpageref{ex}{0}

\begin{ex} %[2D8H40]
Xét tính đúng sai của các phép tính trong bảng nhân 5:
\choiceTF
{\True $5 \times 4 = 20$.}
{\True $5 \times 8 = 40$.}
{$5 \times 7 = 30$.}
{\True $5 \times 7 = 35$.}
\loigiai{
\begin{itemize}
    \item $5 \times 4 = 20$: Đúng.
    \item $5 \times 8 = 40$: Đúng.
    \item $5 \times 7 = 35 \neq 30$. Ý 3 Sai, Ý 4 Đúng.
\end{itemize}
\begin{itemize}
    \item a) \textbf{Đúng}.
    \item b) \textbf{Đúng}.
    \item c) \textbf{Sai}.
    \item d) \textbf{Đúng}.
\end{itemize}
}
\end{ex}

\begin{ex} %[2D8V40]
Một can đựng $5\text{ l}$ nước mắm. Xét tính đúng sai:
\choiceTF
{\True 5 can như thế đựng $25\text{ l}$ nước mắm ($5 \times 5 = 25$).}
{\True 9 can như thế đựng $45\text{ l}$ nước mắm ($5 \times 9 = 45$).}
{10 can như thế đựng $40\text{ l}$ nước mắm.}
{\True 10 can như thế đựng $5 \times 10 = 50\text{ l}$ nước mắm.}
\loigiai{
\begin{itemize}
    \item 5 can: $5 \times 5 = 25\text{ l}$: Đúng.
    \item 9 can: $5 \times 9 = 45\text{ l}$: Đúng.
    \item 10 can: $50\text{ l} \neq 40\text{ l}$. Ý 3 Sai, Ý 4 Đúng.
\end{itemize}
\begin{itemize}
    \item a) \textbf{Đúng}.
    \item b) \textbf{Đúng}.
    \item c) \textbf{Sai}.
    \item d) \textbf{Đúng}.
\end{itemize}
}
\end{ex}

\textbf{Phần 3. Trả lời ngắn}
\setcounterpageref{ex}{0}

\begin{ex} %[2D8H40]
Tính: $5 \times 5$.
\shortans{25}
\loigiai{
\textbf{Đáp số:} $25$\\
$5 \times 5 = 25$.
}
\end{ex}

\begin{ex} %[2D8H40]
Tính: $5 \times 9$.
\shortans{45}
\loigiai{
\textbf{Đáp số:} $45$\\
$5 \times 9 = 45$.
}
\end{ex}

\begin{ex} %[2D8V40]
Tính giá trị biểu thức: $5 \times 8 - 15$.
\shortans{25}
\loigiai{
\textbf{Đáp số:} $25$\\
$40 - 15 = 25$.
}
\end{ex}

\begin{ex} %[2D8V40]
Mỗi gói bánh có 5 cái bánh. Hỏi 6 gói bánh có bao nhiêu cái bánh?
\shortans{30}
\loigiai{
\textbf{Đáp số:} $30$\\
$5 \times 6 = 30$ cái bánh.
}
\end{ex}

\begin{ex} %[2D8V40]
Điền số thích hợp: $5 \times ? = 35$.
\shortans{7}
\loigiai{
\textbf{Đáp số:} $7$\\
$35 : 5 = 7$.
}
\end{ex}

\begin{ex} %[2D8C40]
Một nhóm có 8 bạn, mỗi bạn góp 5 quyển vở ủng hộ vùng cao. Sau đó cả nhóm mua thêm 10 quyển vở nữa. Hỏi nhóm đã ủng hộ tất cả bao nhiêu quyển vở?
\shortans{50}
\loigiai{
\textbf{Đáp số:} $50$\\
$5 \times 8 + 10 = 40 + 10 = 50$ quyển.
}
\end{ex}
